
\documentclass{article}
\usepackage{colortbl}
\usepackage{makecell}
\usepackage{multirow}
\usepackage{supertabular}

\begin{document}

\newcounter{utterance}

\centering \large Interaction Transcript for game `matchit\_ascii', experiment `same\_grid', episode 1 with claude{-}opus{-}4{-}8{-}azure.
\vspace{24pt}

{ \footnotesize  \setcounter{utterance}{1}
\setlength{\tabcolsep}{0pt}
\begin{supertabular}{c@{$\;$}|p{.15\linewidth}@{}p{.15\linewidth}p{.15\linewidth}p{.15\linewidth}p{.15\linewidth}p{.15\linewidth}}
    \# & $\;$A & \multicolumn{4}{c}{Game Master} & $\;\:$B\\
    \hline

    \theutterance \stepcounter{utterance}  
    & & \multicolumn{4}{p{0.6\linewidth}}{
        \cellcolor[rgb]{0.9,0.9,0.9}{
            \makecell[{{p{\linewidth}}}]{
                \texttt{\tiny{[P1$\langle$GM]}}
                \texttt{Osallistut yhteistyöhön perustuvaan arvauspeliin.} \\
\texttt{Tavoitteena on selvittää, ovatko tämä ruudukko ja toinen kuva, jonka vain minä näen, samat.} \\
\texttt{Kuvaile ensin ruudukkosi. Sitten annan oman kuvaukseni ja voimme esittää toisillemme kysymyksiä ruudukoista selvittääksemme, ovatko ne samat. Aloita nyt lyhyt ruudukko{-}kuvauksesi merkillä "KUVAUS:" ja sen jälkeen kuvaus. Älä lisää mitään muuta. Tämä on sinun ruudukkosi:} \\
\texttt{▢ ▢ X ▢ ▢} \\
\texttt{▢ ▢ X ▢ ▢} \\
\texttt{▢ ▢ X ▢ ▢} \\
\texttt{▢ ▢ X ▢ ▢} \\
\texttt{X X X X X} \\
            }
        }
    }
    & \\ \\

    \theutterance \stepcounter{utterance}  
    & \multicolumn{4}{p{0.6\linewidth}}{
        \cellcolor[rgb]{0.9,0.9,0.9}{
            \makecell[{{p{\linewidth}}}]{
                \texttt{\tiny{[P1$\rangle$GM]}}
                \texttt{KUVAUS: Ruudukko on 5×5. Kolmannessa sarakkeessa on pystysuora X{-}viiva ylhäältä alas, ja alimmalla rivillä on vaakasuora X{-}viiva vasemmalta oikealle. Muut ruudut ovat tyhjiä. Yhdessä ne muodostavat ylösalaisin olevan T{-}kirjaimen tai plus{-}merkin alaosan.} \\
            }
        }
    }
    & & \\ \\

    \theutterance \stepcounter{utterance}  
    & & & \multicolumn{2}{p{0.3\linewidth}}{
        \cellcolor[rgb]{0.9,0.9,0.9}{
            \makecell[{{p{\linewidth}}}]{
                \texttt{\tiny{[GM$|$GM]}}
                \texttt{continue} \\
            }
        }
    }
    & & \\ \\

    \theutterance \stepcounter{utterance}  
    & & \multicolumn{4}{p{0.6\linewidth}}{
        \cellcolor[rgb]{0.9,0.9,0.9}{
            \makecell[{{p{\linewidth}}}]{
                \texttt{\tiny{[GM$\rangle$P2]}}
                \texttt{Osallistut yhteistyöhön perustuvaan arvauspeliin.} \\
\texttt{Tavoitteena on selvittää, ovatko tämä ruudukko ja toinen kuva, jonka vain minä näen, samat.} \\
\texttt{Kuvaile ensin ruudukkosi. Sitten annan oman kuvaukseni ja voimme esittää toisillemme kysymyksiä ruudukoista selvittääksemme, ovatko ne samat. Aloita nyt lyhyt ruudukko{-}kuvauksesi merkillä "KUVAUS:" ja sen jälkeen kuvaus. Älä lisää mitään muuta. Tämä on sinun ruudukkosi:} \\
\texttt{▢ ▢ X ▢ ▢} \\
\texttt{▢ ▢ X ▢ ▢} \\
\texttt{▢ ▢ X ▢ ▢} \\
\texttt{▢ ▢ X ▢ ▢} \\
\texttt{X X X X X} \\
            }
        }
    }
    & \\ \\

    \theutterance \stepcounter{utterance}  
    & & & \multicolumn{4}{p{0.6\linewidth}}{
        \cellcolor[rgb]{0.9,0.9,0.9}{
            \makecell[{{p{\linewidth}}}]{
                \texttt{\tiny{[GM$\langle$P2]}}
                \texttt{KUVAUS: Ruudukko on 5×5. Keskimmäinen sarake (kolmas) on täynnä X{-}merkkejä ylhäältä alas, ja alin rivi on kokonaan X{-}merkkejä. Muut ruudut ovat tyhjiä. Muodostuu ristikkäinen kuvio, jossa pystysuora viiva keskellä ja vaakasuora viiva alhaalla, kuten iso "T" ylösalaisin tai plus{-}merkin osa.} \\
            }
        }
    }
     \\ \\

    \theutterance \stepcounter{utterance}  
    & & & \multicolumn{2}{p{0.3\linewidth}}{
        \cellcolor[rgb]{0.9,0.9,0.9}{
            \makecell[{{p{\linewidth}}}]{
                \texttt{\tiny{[GM$|$GM]}}
                \texttt{continue} \\
            }
        }
    }
    & & \\ \\

    \theutterance \stepcounter{utterance}  
    & & \multicolumn{4}{p{0.6\linewidth}}{
        \cellcolor[rgb]{0.9,0.9,0.9}{
            \makecell[{{p{\linewidth}}}]{
                \texttt{\tiny{[GM$\rangle$P2]}}
                \texttt{Tämä on KUVAUKSENIKUVAUS: Ruudukko on 5×5. Kolmannessa sarakkeessa on pystysuora X{-}viiva ylhäältä alas, ja alimmalla rivillä on vaakasuora X{-}viiva vasemmalta oikealle. Muut ruudut ovat tyhjiä. Yhdessä ne muodostavat ylösalaisin olevan T{-}kirjaimen tai plus{-}merkin alaosan.} \\
\texttt{Esitä nyt kysymys selvittääksesi uusia näkökohtia ruudukostani, jotka saattavat poiketa sinun ruudukostasi. Aloita "KYSYMYS:" äläkä lisää mitään muuta.} \\
            }
        }
    }
    & \\ \\

    \theutterance \stepcounter{utterance}  
    & & & \multicolumn{4}{p{0.6\linewidth}}{
        \cellcolor[rgb]{0.9,0.9,0.9}{
            \makecell[{{p{\linewidth}}}]{
                \texttt{\tiny{[GM$\langle$P2]}}
                \texttt{KYSYMYS: Ulottuuko pystysuora X{-}viiva kolmannessa sarakkeessa aivan ylimmältä riviltä (rivi 1) alimmalle riville (rivi 5) asti, eli onko kolmannessa sarakkeessa X jokaisella viidellä rivillä?} \\
            }
        }
    }
     \\ \\

    \theutterance \stepcounter{utterance}  
    & & & \multicolumn{2}{p{0.3\linewidth}}{
        \cellcolor[rgb]{0.9,0.9,0.9}{
            \makecell[{{p{\linewidth}}}]{
                \texttt{\tiny{[GM$|$GM]}}
                \texttt{continue} \\
            }
        }
    }
    & & \\ \\

    \theutterance \stepcounter{utterance}  
    & & \multicolumn{4}{p{0.6\linewidth}}{
        \cellcolor[rgb]{0.9,0.9,0.9}{
            \makecell[{{p{\linewidth}}}]{
                \texttt{\tiny{[P1$\langle$GM]}}
                \texttt{Tämä on KUVAUKSENIKUVAUS: Ruudukko on 5×5. Keskimmäinen sarake (kolmas) on täynnä X{-}merkkejä ylhäältä alas, ja alin rivi on kokonaan X{-}merkkejä. Muut ruudut ovat tyhjiä. Muodostuu ristikkäinen kuvio, jossa pystysuora viiva keskellä ja vaakasuora viiva alhaalla, kuten iso "T" ylösalaisin tai plus{-}merkin osa.} \\
\texttt{KYSYMYS: Ulottuuko pystysuora X{-}viiva kolmannessa sarakkeessa aivan ylimmältä riviltä (rivi 1) alimmalle riville (rivi 5) asti, eli onko kolmannessa sarakkeessa X jokaisella viidellä rivillä?Aloita vastauksesi merkillä "VASTAUS:" äläkä lisää mitään muuta.} \\
            }
        }
    }
    & \\ \\

    \theutterance \stepcounter{utterance}  
    & \multicolumn{4}{p{0.6\linewidth}}{
        \cellcolor[rgb]{0.9,0.9,0.9}{
            \makecell[{{p{\linewidth}}}]{
                \texttt{\tiny{[P1$\rangle$GM]}}
                \texttt{VASTAUS: Kyllä, kolmannessa sarakkeessa on X jokaisella viidellä rivillä ylhäältä alas. Ruudukkomme vaikuttavat identtisiltä.} \\
            }
        }
    }
    & & \\ \\

    \theutterance \stepcounter{utterance}  
    & & & \multicolumn{2}{p{0.3\linewidth}}{
        \cellcolor[rgb]{0.9,0.9,0.9}{
            \makecell[{{p{\linewidth}}}]{
                \texttt{\tiny{[GM$|$GM]}}
                \texttt{continue} \\
            }
        }
    }
    & & \\ \\

    \theutterance \stepcounter{utterance}  
    & & \multicolumn{4}{p{0.6\linewidth}}{
        \cellcolor[rgb]{0.9,0.9,0.9}{
            \makecell[{{p{\linewidth}}}]{
                \texttt{\tiny{[P1$\langle$GM]}}
                \texttt{Esitä nyt kysymys selvittääksesi uusia näkökohtia ruudukostani, jotka saattavat poiketa sinun ruudukostasi. Aloita "KYSYMYS:" äläkä lisää mitään muuta.} \\
            }
        }
    }
    & \\ \\

    \theutterance \stepcounter{utterance}  
    & \multicolumn{4}{p{0.6\linewidth}}{
        \cellcolor[rgb]{0.9,0.9,0.9}{
            \makecell[{{p{\linewidth}}}]{
                \texttt{\tiny{[P1$\rangle$GM]}}
                \texttt{KYSYMYS: Onko alimmalla rivillä (rivi 5) X jokaisessa viidessä sarakkeessa vasemmalta oikealle, eli onko koko alarivi täynnä X{-}merkkejä?} \\
            }
        }
    }
    & & \\ \\

    \theutterance \stepcounter{utterance}  
    & & & \multicolumn{2}{p{0.3\linewidth}}{
        \cellcolor[rgb]{0.9,0.9,0.9}{
            \makecell[{{p{\linewidth}}}]{
                \texttt{\tiny{[GM$|$GM]}}
                \texttt{continue} \\
            }
        }
    }
    & & \\ \\

    \theutterance \stepcounter{utterance}  
    & & \multicolumn{4}{p{0.6\linewidth}}{
        \cellcolor[rgb]{0.9,0.9,0.9}{
            \makecell[{{p{\linewidth}}}]{
                \texttt{\tiny{[GM$\rangle$P2]}}
                \texttt{VASTAUS: Kyllä, kolmannessa sarakkeessa on X jokaisella viidellä rivillä ylhäältä alas. Ruudukkomme vaikuttavat identtisiltä.} \\
\texttt{KYSYMYS: Onko alimmalla rivillä (rivi 5) X jokaisessa viidessä sarakkeessa vasemmalta oikealle, eli onko koko alarivi täynnä X{-}merkkejä?Aloita vastauksesi merkillä "VASTAUS:" äläkä lisää mitään muuta.} \\
            }
        }
    }
    & \\ \\

    \theutterance \stepcounter{utterance}  
    & & & \multicolumn{4}{p{0.6\linewidth}}{
        \cellcolor[rgb]{0.9,0.9,0.9}{
            \makecell[{{p{\linewidth}}}]{
                \texttt{\tiny{[GM$\langle$P2]}}
                \texttt{VASTAUS: Kyllä, alimmalla rivillä on X jokaisessa viidessä sarakkeessa vasemmalta oikealle. Ruudukkomme ovat identtiset.} \\
            }
        }
    }
     \\ \\

    \theutterance \stepcounter{utterance}  
    & & & \multicolumn{2}{p{0.3\linewidth}}{
        \cellcolor[rgb]{0.9,0.9,0.9}{
            \makecell[{{p{\linewidth}}}]{
                \texttt{\tiny{[GM$|$GM]}}
                \texttt{continue} \\
            }
        }
    }
    & & \\ \\

    \theutterance \stepcounter{utterance}  
    & & \multicolumn{4}{p{0.6\linewidth}}{
        \cellcolor[rgb]{0.9,0.9,0.9}{
            \makecell[{{p{\linewidth}}}]{
                \texttt{\tiny{[GM$\rangle$P2]}}
                \texttt{Esitä nyt kysymys selvittääksesi uusia näkökohtia ruudukostani, jotka saattavat poiketa sinun ruudukostasi. Aloita "KYSYMYS:" äläkä lisää mitään muuta.} \\
            }
        }
    }
    & \\ \\

    \theutterance \stepcounter{utterance}  
    & & & \multicolumn{4}{p{0.6\linewidth}}{
        \cellcolor[rgb]{0.9,0.9,0.9}{
            \makecell[{{p{\linewidth}}}]{
                \texttt{\tiny{[GM$\langle$P2]}}
                \texttt{KYSYMYS: Onko kaikissa muissa ruuduissa paitsi kolmannessa sarakkeessa ja alimmalla rivillä tyhjää, eli onko sarakkeissa 1, 2, 4 ja 5 riveillä 1–4 pelkkiä tyhjiä ruutuja?} \\
            }
        }
    }
     \\ \\

    \theutterance \stepcounter{utterance}  
    & & & \multicolumn{2}{p{0.3\linewidth}}{
        \cellcolor[rgb]{0.9,0.9,0.9}{
            \makecell[{{p{\linewidth}}}]{
                \texttt{\tiny{[GM$|$GM]}}
                \texttt{continue} \\
            }
        }
    }
    & & \\ \\

    \theutterance \stepcounter{utterance}  
    & & \multicolumn{4}{p{0.6\linewidth}}{
        \cellcolor[rgb]{0.9,0.9,0.9}{
            \makecell[{{p{\linewidth}}}]{
                \texttt{\tiny{[P1$\langle$GM]}}
                \texttt{VASTAUS: Kyllä, alimmalla rivillä on X jokaisessa viidessä sarakkeessa vasemmalta oikealle. Ruudukkomme ovat identtiset.} \\
\texttt{KYSYMYS: Onko kaikissa muissa ruuduissa paitsi kolmannessa sarakkeessa ja alimmalla rivillä tyhjää, eli onko sarakkeissa 1, 2, 4 ja 5 riveillä 1–4 pelkkiä tyhjiä ruutuja?Aloita vastauksesi merkillä "VASTAUS:" äläkä lisää mitään muuta.} \\
            }
        }
    }
    & \\ \\

    \theutterance \stepcounter{utterance}  
    & \multicolumn{4}{p{0.6\linewidth}}{
        \cellcolor[rgb]{0.9,0.9,0.9}{
            \makecell[{{p{\linewidth}}}]{
                \texttt{\tiny{[P1$\rangle$GM]}}
                \texttt{VASTAUS: Kyllä, sarakkeissa 1, 2, 4 ja 5 riveillä 1–4 on pelkkiä tyhjiä ruutuja. Ruudukkomme ovat identtiset.} \\
            }
        }
    }
    & & \\ \\

    \theutterance \stepcounter{utterance}  
    & & & \multicolumn{2}{p{0.3\linewidth}}{
        \cellcolor[rgb]{0.9,0.9,0.9}{
            \makecell[{{p{\linewidth}}}]{
                \texttt{\tiny{[GM$|$GM]}}
                \texttt{continue} \\
            }
        }
    }
    & & \\ \\

    \theutterance \stepcounter{utterance}  
    & & \multicolumn{4}{p{0.6\linewidth}}{
        \cellcolor[rgb]{0.9,0.9,0.9}{
            \makecell[{{p{\linewidth}}}]{
                \texttt{\tiny{[P1$\langle$GM]}}
                \texttt{Esitä nyt kysymys selvittääksesi uusia näkökohtia ruudukostani, jotka saattavat poiketa sinun ruudukostasi. Aloita "KYSYMYS:" äläkä lisää mitään muuta.} \\
            }
        }
    }
    & \\ \\

    \theutterance \stepcounter{utterance}  
    & \multicolumn{4}{p{0.6\linewidth}}{
        \cellcolor[rgb]{0.9,0.9,0.9}{
            \makecell[{{p{\linewidth}}}]{
                \texttt{\tiny{[P1$\rangle$GM]}}
                \texttt{KYSYMYS: Olemme nyt vahvistaneet jokaisen ruudun tilan. Onko ruudukossasi yhteensä tasan 9 X{-}merkkiä (4 pystyviivassa riveillä 1–4 plus 5 alarivillä)?} \\
            }
        }
    }
    & & \\ \\

    \theutterance \stepcounter{utterance}  
    & & & \multicolumn{2}{p{0.3\linewidth}}{
        \cellcolor[rgb]{0.9,0.9,0.9}{
            \makecell[{{p{\linewidth}}}]{
                \texttt{\tiny{[GM$|$GM]}}
                \texttt{continue} \\
            }
        }
    }
    & & \\ \\

    \theutterance \stepcounter{utterance}  
    & & \multicolumn{4}{p{0.6\linewidth}}{
        \cellcolor[rgb]{0.9,0.9,0.9}{
            \makecell[{{p{\linewidth}}}]{
                \texttt{\tiny{[GM$\rangle$P2]}}
                \texttt{VASTAUS: Kyllä, sarakkeissa 1, 2, 4 ja 5 riveillä 1–4 on pelkkiä tyhjiä ruutuja. Ruudukkomme ovat identtiset.} \\
\texttt{KYSYMYS: Olemme nyt vahvistaneet jokaisen ruudun tilan. Onko ruudukossasi yhteensä tasan 9 X{-}merkkiä (4 pystyviivassa riveillä 1–4 plus 5 alarivillä)?Aloita vastauksesi merkillä "VASTAUS:" äläkä lisää mitään muuta.} \\
            }
        }
    }
    & \\ \\

    \theutterance \stepcounter{utterance}  
    & & & \multicolumn{4}{p{0.6\linewidth}}{
        \cellcolor[rgb]{0.9,0.9,0.9}{
            \makecell[{{p{\linewidth}}}]{
                \texttt{\tiny{[GM$\langle$P2]}}
                \texttt{VASTAUS: Kyllä, ruudukossani on tasan 9 X{-}merkkiä (4 pystyviivassa riveillä 1–4 plus 5 alarivillä). Ruudukkomme ovat identtiset.} \\
            }
        }
    }
     \\ \\

    \theutterance \stepcounter{utterance}  
    & & & \multicolumn{2}{p{0.3\linewidth}}{
        \cellcolor[rgb]{0.9,0.9,0.9}{
            \makecell[{{p{\linewidth}}}]{
                \texttt{\tiny{[GM$|$GM]}}
                \texttt{continue} \\
            }
        }
    }
    & & \\ \\

    \theutterance \stepcounter{utterance}  
    & & \multicolumn{4}{p{0.6\linewidth}}{
        \cellcolor[rgb]{0.9,0.9,0.9}{
            \makecell[{{p{\linewidth}}}]{
                \texttt{\tiny{[GM$\rangle$P2]}}
                \texttt{Esitä nyt kysymys selvittääksesi uusia näkökohtia ruudukostani, jotka saattavat poiketa sinun ruudukostasi. Aloita "KYSYMYS:" äläkä lisää mitään muuta.} \\
            }
        }
    }
    & \\ \\

    \theutterance \stepcounter{utterance}  
    & & & \multicolumn{4}{p{0.6\linewidth}}{
        \cellcolor[rgb]{0.9,0.9,0.9}{
            \makecell[{{p{\linewidth}}}]{
                \texttt{\tiny{[GM$\langle$P2]}}
                \texttt{KYSYMYS: Olemme nyt käyneet läpi kaikki ruudut ja X{-}merkkien lukumäärän. Onko mielestäsi vielä jokin yksityiskohta, jota emme ole vahvistaneet, vai voimmeko todeta ruudukot varmasti identtisiksi?} \\
            }
        }
    }
     \\ \\

    \theutterance \stepcounter{utterance}  
    & & & \multicolumn{2}{p{0.3\linewidth}}{
        \cellcolor[rgb]{0.9,0.9,0.9}{
            \makecell[{{p{\linewidth}}}]{
                \texttt{\tiny{[GM$|$GM]}}
                \texttt{continue} \\
            }
        }
    }
    & & \\ \\

    \theutterance \stepcounter{utterance}  
    & & \multicolumn{4}{p{0.6\linewidth}}{
        \cellcolor[rgb]{0.9,0.9,0.9}{
            \makecell[{{p{\linewidth}}}]{
                \texttt{\tiny{[P1$\langle$GM]}}
                \texttt{VASTAUS: Kyllä, ruudukossani on tasan 9 X{-}merkkiä (4 pystyviivassa riveillä 1–4 plus 5 alarivillä). Ruudukkomme ovat identtiset.} \\
\texttt{KYSYMYS: Olemme nyt käyneet läpi kaikki ruudut ja X{-}merkkien lukumäärän. Onko mielestäsi vielä jokin yksityiskohta, jota emme ole vahvistaneet, vai voimmeko todeta ruudukot varmasti identtisiksi?Aloita vastauksesi merkillä "VASTAUS:" äläkä lisää mitään muuta.} \\
            }
        }
    }
    & \\ \\

    \theutterance \stepcounter{utterance}  
    & \multicolumn{4}{p{0.6\linewidth}}{
        \cellcolor[rgb]{0.9,0.9,0.9}{
            \makecell[{{p{\linewidth}}}]{
                \texttt{\tiny{[P1$\rangle$GM]}}
                \texttt{VASTAUS: Olemme nyt vahvistaneet kaikki yksityiskohdat: kolmannen sarakkeen pystyviivan, alarivin vaakaviivan, muiden ruutujen tyhjyyden sekä X{-}merkkien kokonaismäärän. Voimme todeta ruudukot varmasti identtisiksi.} \\
            }
        }
    }
    & & \\ \\

    \theutterance \stepcounter{utterance}  
    & & & \multicolumn{2}{p{0.3\linewidth}}{
        \cellcolor[rgb]{0.9,0.9,0.9}{
            \makecell[{{p{\linewidth}}}]{
                \texttt{\tiny{[GM$|$GM]}}
                \texttt{continue} \\
            }
        }
    }
    & & \\ \\

    \theutterance \stepcounter{utterance}  
    & & \multicolumn{4}{p{0.6\linewidth}}{
        \cellcolor[rgb]{0.9,0.9,0.9}{
            \makecell[{{p{\linewidth}}}]{
                \texttt{\tiny{[P1$\langle$GM]}}
                \texttt{Esitä nyt kysymys selvittääksesi uusia näkökohtia ruudukostani, jotka saattavat poiketa sinun ruudukostasi. Aloita "KYSYMYS:" äläkä lisää mitään muuta.} \\
            }
        }
    }
    & \\ \\

    \theutterance \stepcounter{utterance}  
    & \multicolumn{4}{p{0.6\linewidth}}{
        \cellcolor[rgb]{0.9,0.9,0.9}{
            \makecell[{{p{\linewidth}}}]{
                \texttt{\tiny{[P1$\rangle$GM]}}
                \texttt{KYSYMYS: Vahvistetaan vielä yksi risteyskohta: onko kolmannen sarakkeen ja alarivin leikkauskohdassa eli ruudussa (rivi 5, sarake 3) X{-}merkki?} \\
            }
        }
    }
    & & \\ \\

    \theutterance \stepcounter{utterance}  
    & & & \multicolumn{2}{p{0.3\linewidth}}{
        \cellcolor[rgb]{0.9,0.9,0.9}{
            \makecell[{{p{\linewidth}}}]{
                \texttt{\tiny{[GM$|$GM]}}
                \texttt{continue} \\
            }
        }
    }
    & & \\ \\

    \theutterance \stepcounter{utterance}  
    & & \multicolumn{4}{p{0.6\linewidth}}{
        \cellcolor[rgb]{0.9,0.9,0.9}{
            \makecell[{{p{\linewidth}}}]{
                \texttt{\tiny{[GM$\rangle$P2]}}
                \texttt{VASTAUS: Olemme nyt vahvistaneet kaikki yksityiskohdat: kolmannen sarakkeen pystyviivan, alarivin vaakaviivan, muiden ruutujen tyhjyyden sekä X{-}merkkien kokonaismäärän. Voimme todeta ruudukot varmasti identtisiksi.} \\
\texttt{KYSYMYS: Vahvistetaan vielä yksi risteyskohta: onko kolmannen sarakkeen ja alarivin leikkauskohdassa eli ruudussa (rivi 5, sarake 3) X{-}merkki?Aloita vastauksesi merkillä "VASTAUS:" äläkä lisää mitään muuta.} \\
            }
        }
    }
    & \\ \\

    \theutterance \stepcounter{utterance}  
    & & & \multicolumn{4}{p{0.6\linewidth}}{
        \cellcolor[rgb]{0.9,0.9,0.9}{
            \makecell[{{p{\linewidth}}}]{
                \texttt{\tiny{[GM$\langle$P2]}}
                \texttt{VASTAUS: Kyllä, ruudussa (rivi 5, sarake 3) on X{-}merkki, sillä se on osa sekä pystyviivaa että alariviä. Ruudukkomme ovat identtiset.} \\
            }
        }
    }
     \\ \\

    \theutterance \stepcounter{utterance}  
    & & & \multicolumn{2}{p{0.3\linewidth}}{
        \cellcolor[rgb]{0.9,0.9,0.9}{
            \makecell[{{p{\linewidth}}}]{
                \texttt{\tiny{[GM$|$GM]}}
                \texttt{continue} \\
            }
        }
    }
    & & \\ \\

    \theutterance \stepcounter{utterance}  
    & & \multicolumn{4}{p{0.6\linewidth}}{
        \cellcolor[rgb]{0.9,0.9,0.9}{
            \makecell[{{p{\linewidth}}}]{
                \texttt{\tiny{[GM$\rangle$P2]}}
                \texttt{Tee nyt päätös. Mitä mieltä olet: ovatko ruudukkosi ja toinen kuvattu ruudukko samat? Kirjoita "PÄÄTÖS: sama ruudukko", jos uskot niiden olevan sama ruudukko, tai "PÄÄTÖS: eri ruudukot", jos uskot niiden olevan eri ruudukot. Älä lisää mitään muuta.} \\
            }
        }
    }
    & \\ \\

    \theutterance \stepcounter{utterance}  
    & & & \multicolumn{4}{p{0.6\linewidth}}{
        \cellcolor[rgb]{0.9,0.9,0.9}{
            \makecell[{{p{\linewidth}}}]{
                \texttt{\tiny{[GM$\langle$P2]}}
                \texttt{PÄÄTÖS: sama ruudukko} \\
            }
        }
    }
     \\ \\

    \theutterance \stepcounter{utterance}  
    & & & \multicolumn{2}{p{0.3\linewidth}}{
        \cellcolor[rgb]{0.9,0.9,0.9}{
            \makecell[{{p{\linewidth}}}]{
                \texttt{\tiny{[GM$|$GM]}}
                \texttt{continue} \\
            }
        }
    }
    & & \\ \\

    \theutterance \stepcounter{utterance}  
    & & & \multicolumn{2}{p{0.3\linewidth}}{
        \cellcolor[rgb]{0.9,0.9,0.9}{
            \makecell[{{p{\linewidth}}}]{
                \texttt{\tiny{[GM$|$GM]}}
                \texttt{success} \\
            }
        }
    }
    & & \\ \\

    \theutterance \stepcounter{utterance}  
    & & \multicolumn{4}{p{0.6\linewidth}}{
        \cellcolor[rgb]{0.9,0.9,0.9}{
            \makecell[{{p{\linewidth}}}]{
                \texttt{\tiny{[P1$\langle$GM]}}
                \texttt{VASTAUS: Kyllä, ruudussa (rivi 5, sarake 3) on X{-}merkki, sillä se on osa sekä pystyviivaa että alariviä. Ruudukkomme ovat identtiset.} \\
\texttt{Tee nyt päätös. Mitä mieltä olet: ovatko ruudukkosi ja toinen kuvattu ruudukko samat? Kirjoita "PÄÄTÖS: sama ruudukko", jos uskot niiden olevan sama ruudukko, tai "PÄÄTÖS: eri ruudukot", jos uskot niiden olevan eri ruudukot. Älä lisää mitään muuta.} \\
            }
        }
    }
    & \\ \\

    \theutterance \stepcounter{utterance}  
    & \multicolumn{4}{p{0.6\linewidth}}{
        \cellcolor[rgb]{0.9,0.9,0.9}{
            \makecell[{{p{\linewidth}}}]{
                \texttt{\tiny{[P1$\rangle$GM]}}
                \texttt{PÄÄTÖS: sama ruudukko} \\
            }
        }
    }
    & & \\ \\

    \theutterance \stepcounter{utterance}  
    & & & \multicolumn{2}{p{0.3\linewidth}}{
        \cellcolor[rgb]{0.9,0.9,0.9}{
            \makecell[{{p{\linewidth}}}]{
                \texttt{\tiny{[GM$|$GM]}}
                \texttt{continue} \\
            }
        }
    }
    & & \\ \\

    \theutterance \stepcounter{utterance}  
    & & & \multicolumn{2}{p{0.3\linewidth}}{
        \cellcolor[rgb]{0.9,0.9,0.9}{
            \makecell[{{p{\linewidth}}}]{
                \texttt{\tiny{[GM$|$GM]}}
                \texttt{success} \\
            }
        }
    }
    & & \\ \\

\end{supertabular}
}

\end{document}
